% =============================================================================
% CAPÍTULO 2 — ARQUITECTURA DE ACCESO Y CONCEPTOS BÁSICOS
% =============================================================================
\chapter{Arquitectura de Acceso y Conceptos Básicos}
\label{ch:acceso}
\resumenCapitulo{Este capítulo describe el modelo de perfiles y permisos de EthosHub,
los requisitos mínimos del entorno para utilizar la plataforma y los procedimientos de
acceso inicial: registro, inicio de sesión y recuperación de credenciales.}

% -----------------------------------------------------------------------------
\section{Matriz de Perfiles y Roles (RBAC)}
\label{sec:rbac}

EthosHub implementa un control de acceso basado en roles (\textit{Role-Based Access
Control}, RBAC). Al registrarse, cada persona usuaria adopta un perfil que determina
las funciones a las que puede acceder. Los perfiles son mutuamente excluyentes: una
cuenta opera como Profesional, como Reclutador o como Administrador.

\vspace{0.2cm}
\noindent
\renewcommand{\arraystretch}{1.4}
\fontsize{9}{11}\selectfont\sffamily
\begin{tabularx}{\textwidth}{|>{\bfseries\color{colorPrimario}\raggedright\arraybackslash}X|C{2.6cm}|C{2.6cm}|C{2.6cm}|}
\hline
\rowcolor{headerTabla}
\sffamily\textbf{Función / Capacidad} & \sffamily\textbf{Profesional} & \sffamily\textbf{Reclutador} & \sffamily\textbf{Administrador} \\
\hline
Gestionar portafolio propio (proyectos, experiencia, educación, habilidades) & Sí & No & No \\
\hline
\rowcolor{grisTecnico}
Generar currículum (CV Studio) & Sí & No & No \\
\hline
Buscar y filtrar talento & No & Sí & Sí \\
\hline
\rowcolor{grisTecnico}
Gestionar \textit{pipeline} de candidatos & No & Sí & No \\
\hline
Contactar por chat & Sí & Sí & No \\
\hline
\rowcolor{grisTecnico}
Moderar contenido y reportes & No & No & Sí \\
\hline
Gestionar cuentas, roles y permisos & No & No & Sí \\
\hline
\rowcolor{grisTecnico}
Administrar el repositorio global de habilidades & No & No & Sí \\
\hline
Consultar métricas globales del sistema & No & No & Sí \\
\hline
\rowcolor{grisTecnico}
Configurar su propia cuenta (perfil, seguridad, datos) & Sí & Sí & Sí \\
\hline
\end{tabularx}

\vspace{0.3cm}
\begin{AvisoNota}
Si necesita operar con un perfil distinto al de su cuenta actual (por ejemplo, pasar de
Profesional a Reclutador), deberá solicitarlo al equipo de administración; el cambio de
rol no puede realizarse de forma autónoma.
\end{AvisoNota}

% -----------------------------------------------------------------------------
\section{Requisitos Mínimos del Entorno}
\label{sec:requisitos}

EthosHub es una aplicación web; no requiere instalación local. Para una experiencia
óptima se recomienda el siguiente entorno.

\begin{TablaInfo}{Componente}{Requisito recomendado}
\textbf{Navegador web} & Google Chrome, Microsoft Edge, Mozilla Firefox o Safari, en
su versión estable más reciente (o las dos versiones anteriores). \\ \hline
\rowcolor{grisTecnico}
\textbf{Conectividad} & Conexión a Internet estable de al menos 2~Mbps para la carga de
capturas, documentos y multimedia. \\ \hline
\textbf{Resolución} & 1280\,$\times$\,720 píxeles o superior en escritorio. La
plataforma también ofrece diseño adaptable para teléfonos y tabletas. \\ \hline
\rowcolor{grisTecnico}
\textbf{Cookies y JavaScript} & Habilitados, requisito indispensable para la sesión y
las funciones interactivas. \\ \hline
\textbf{Cuenta de correo} & Una dirección de correo válida y accesible, necesaria para
la verificación y la recuperación de credenciales. \\ \hline
\end{TablaInfo}

\begin{AvisoPrecaucion}
El uso de navegadores obsoletos o con extensiones que bloqueen JavaScript puede impedir
el funcionamiento correcto de formularios, mapas interactivos y carga de archivos.
\end{AvisoPrecaucion}

% -----------------------------------------------------------------------------
\section{Gestión de Acceso Inicial}
\label{sec:acceso-inicial}

Antes de utilizar cualquier función, la persona usuaria debe disponer de una cuenta y
haber iniciado sesión. A continuación se describen los tres procedimientos de acceso.

\subsection{Registro de Cuenta Nueva}
\label{subsec:registro}

\begin{enumerate}[noitemsep]
    \item Abra el navegador e ingrese a la dirección de la plataforma:
    \comando{https://ethoshub.com}.
    \item En la pantalla de bienvenida, pulse el botón \boton{Registrarse} (o
    \boton{Crear cuenta}).
    \item Complete el formulario de registro con su \campo{Nombre}, \campo{Apellido},
    \campo{Correo electrónico} y \campo{Contraseña}. La contraseña debe cumplir los
    requisitos mínimos de seguridad indicados en pantalla.
    \item Seleccione el tipo de perfil con el que desea operar (\campo{Profesional} o
    \campo{Reclutador}).
    \item Pulse \boton{Crear cuenta}. El sistema enviará un código o enlace de
    verificación a su correo electrónico.
    \item Abra el mensaje recibido e introduzca el código de verificación, o pulse el
    enlace, para activar la cuenta.
\end{enumerate}

\resultado{Resultado esperado:} la cuenta queda activada y el sistema lo dirige al
espacio de trabajo correspondiente a su perfil.

\subsection{Inicio de Sesión Autónomo}
\label{subsec:login}

\begin{enumerate}[noitemsep]
    \item Ingrese a \comando{https://ethoshub.com} y pulse \boton{Iniciar sesión}.
    \item Introduzca el \campo{Correo electrónico} y la \campo{Contraseña} asociados a
    su cuenta.
    \item Pulse \boton{Entrar}. Si las credenciales son correctas, accederá a su panel
    principal.
\end{enumerate}

\begin{AvisoNota}
Si activa la opción \campo{Recordarme}, la sesión se mantendrá iniciada en ese
navegador hasta que cierre sesión manualmente. Evite usar esta opción en equipos
compartidos.
\end{AvisoNota}

\subsection{Recuperación de Credenciales de Acceso}
\label{subsec:recuperacion}

Si olvidó su contraseña, puede restablecerla sin asistencia:

\begin{enumerate}[noitemsep]
    \item En la pantalla de inicio de sesión, pulse el enlace
    \boton{¿Olvidó su contraseña?}.
    \item Introduzca el \campo{Correo electrónico} asociado a la cuenta y pulse
    \boton{Enviar código}.
    \item Revise su bandeja de entrada e introduzca el código de verificación recibido.
    \item Defina una \campo{Nueva contraseña}, confírmela y pulse
    \boton{Actualizar contraseña}.
\end{enumerate}

\resultado{Resultado esperado:} la contraseña se actualiza y puede iniciar sesión con
las nuevas credenciales.

\begin{AvisoPrecaucion}
El código de verificación tiene una vigencia limitada. Si caduca, repita el
procedimiento para solicitar uno nuevo.
\end{AvisoPrecaucion}
